\documentclass[11pt,a4paper]{article}
\usepackage[utf8]{inputenc}\usepackage[T1]{fontenc}\usepackage[italian,english]{babel}
\usepackage{amsmath,amssymb,amsthm}\usepackage[margin=2.2cm]{geometry}\usepackage{longtable,booktabs,array,calc}
\usepackage{listings,xcolor}\usepackage{hyperref}\usepackage{url}
\providecommand{\tightlist}{\setlength{\itemsep}{0pt}\setlength{\parskip}{0pt}}
\lstset{basicstyle=\ttfamily\scriptsize,breaklines=true,frame=single,captionpos=b,columns=fullflexible,keepspaces=true,inputencoding=utf8,extendedchars=true}
\title{Proof Pursuit --- Problem 3, part 2\\ \large prove, decisioni degli agenti, codice verificato, letteratura}
\author{Team Proof Pursuit (Researcher: T. Tumini; Referee: gabundos; Creative: M. Renzi; gio3r3gio)}
\date{\today}
\begin{document}\maketitle\tableofcontents
\section*{Come leggere questo rapporto}
Per ogni cella: la richiesta ufficiale, lo stato dichiarato dal team (mai ``risolto'' senza revisione umana), la consegna
con la prova, poi la traccia delle decisioni degli agenti (Researcher: famiglia e motivo; Referee: verdetto ed errore
fatale; Creative: quando e perché è intervenuto) e il codice su cui il tentativo poggia con l'esito della riesecuzione
fidata. DIMOSTRATO, CITATO, VERIFICATO SU INTERVALLO FINITO e CONGETTURA restano distinti. L'architettura del sistema
è descritta nel README del repository.
\section{Problem 3 — Bulgarian solitaire}
\subsection{Cella 2 (2 punti) — stato: parziale}
\subparagraph{\texorpdfstring{Parte 2 (C2) --- \(D_B\) at triangular
\(n\)}{Parte 2 (C2) --- D\_B at triangular n}}\label{parte-2-c2-d_b-at-triangular-n}

\textbf{Punteggio:} 2 points · \textbf{Valutazione:} Judged

Next, how long the process can take to reach a cycle, starting with the
triangular numbers. Determine \(D_B(T_k)\) for every \(k\), with proof
of both bounds.

\textbf{Enunciato protetto dato agli agenti (state.json).} \texttt{Determine D\_B(T\_k) for every k, with proof of both bounds (exact value as a formula in k).}

\textbf{Evidenza registrata in STATUS.md.} loop p3\_c2: bound inferiore dimostrato con estremale esplicita (lemma di moto delle celle), verifica esaustiva k≤9; bound superiore NON dimostrato (Referee REJECT sul tutto, lemma accettato); consegna parziale in inglese

\subsubsection{Consegna (prova)}
\paragraph{Problem 3 --- Part 2 --- submission
(PARTIAL)}\label{problem-3-part-2-submission-partial}

\textbf{Declared status: PARTIAL.} Lower bound proved in full with an
explicit extremal partition; the upper bound is NOT proved here
(verified exhaustively for small k only, and known in the literature:
Igusa 1985, Etienne 1991, whose proofs we could not access in time). The
automatic Referee rejected the attempt as a solution of the whole cell
precisely for this reason (verdict REJECT: ``Theorem B (upper bound
D\_B(T\_k) ≤ k\^{}2−k for all k) is not proved; the submission itself
marks it `NOT PROVED HERE' and only cites Igusa/Etienne (statement not
reproduced, sources not read) plus an exhaustive check for
sm\ldots{}''), and accepted the lower-bound part as a lemma candidate.

\subparagraph{1. Result and scope}\label{result-and-scope}

Complete proof that D\_B(T\_k) ≥ k\^{}2 − k for all k ≥ 1 with the
explicit extremal γ\_k = (k−1,k−1,k−2,\ldots,2,1,1) (Theorem A), resting
on a new exact cell-motion lemma for B (Lemma 1) that is reusable for
the other cells; exhaustive exact verification that D\_B(T\_k) = k\^{}2
− k for k ≤ 9. The general upper bound D\_B(T\_k) ≤ k\^{}2 − k is not
proved.

\subparagraph{2. Proof}\label{proof}

\subparagraph{Conventions}\label{conventions}

Draw a partition \(\lambda=(\lambda_1\ge\dots\ge\lambda_s)\) as a
diagram whose \textbf{column} \(i\) (\(1\le i\le s\)) consists of the
cells \((i,h)\), \(1\le h\le \lambda_i\) (\(h\) = height). The
\emph{diagonal} of a cell is \(d(i,h)=i+h-1\). Diagonal \(d\) contains
the \(d\) \emph{slots} \((p,\,d+1-p)\), \(p=1,\dots,d\); a set of cells
is the diagram of a partition iff the column heights are weakly
decreasing. \(\delta_k\) is the diagram whose diagonals \(1,\dots,k\)
are full and all others empty. Throughout \(n=T_k\).

\subparagraph{\texorpdfstring{Lemma 1 (exact cell rule for
\(B\))}{Lemma 1 (exact cell rule for B)}}\label{lemma-1-exact-cell-rule-for-b}

Let \(\lambda\) have \(s\) columns. Then \(B(\lambda)\) is the diagram
obtained from \(\lambda\) by moving every cell as follows:

\begin{enumerate}
\def\labelenumi{\arabic{enumi}.}
\tightlist
\item
  \((i,1)\mapsto(1,i)\) for \(1\le i\le s\);
\item
  \((i,h)\mapsto(i+1,h-1)\) if \(2\le h\le s+1\);
\item
  \((i,h)\mapsto(i,h-1)\) if \(h\ge s+2\).
\end{enumerate}

In particular cells of type 1--2 stay on their diagonal (slot
\(p\mapsto p+1\), and slot \(d\mapsto 1\): a cyclic rotation of the
diagonal), while cells of type 3 move from diagonal \(d\) to diagonal
\(d-1\) (same column).

\emph{Proof.} By definition the columns of \(B(\lambda)\) are the
multiset \(\{s\}\cup\{\lambda_i-1:\lambda_i\ge 2\}\). Consider first the
``unsorted'' diagram \(U\): column 1 of height \(s\), column \(i+1\) of
height \(\lambda_i-1\) (\(1\le i\le s\)). Rules 1 and 2 applied to all
cells with \(h\le s+1\) produce exactly the cells of \(U\) of height
\(\le s+1\) (the new column, and the old columns shifted right and
lowered), and rule 2 without the restriction \(h\le s+1\) would produce
all of \(U\). Let
\(q=\#\{i:\lambda_i-1\ge s+1\}=\#\{i:\lambda_i\ge s+2\}\); because
\(\lambda\) is sorted, these are \(i=1,\dots,q\), so in \(U\) exactly
the columns \(2,\dots,q+1\) have height \(>s+1\) (column 1 has height
\(s\), columns \(i+1\) with \(i>q\) have height \(\le s\)). Now compare
\(U\) with the diagram \(V\) obtained from \(U\) by moving every cell of
height \(\ge s+1\) in columns \(2,\dots,q+1\) one column to the left.
Column heights of \(V\): column 1 has the \(s\) cells of the new column
plus the cells of heights \(s+1,\dots,\lambda_1-1\) received from column
2, hence height \(\lambda_1-1\); column \(i\) for \(2\le i\le q\) keeps
heights \(1,\dots,s\) (from old column \(i-1\), whose height
\(\lambda_{i-1}-1\ge s+1\)), loses heights \(s+1,\dots,\lambda_{i-1}-1\)
and receives heights \(s+1,\dots,\lambda_i-1\) from column \(i+1\),
hence height \(\lambda_i-1\); column \(q+1\) keeps heights \(1,\dots,s\)
and loses the rest, hence height \(s\); columns \(i+1>q+1\) are
unchanged with height \(\lambda_i-1\le s\). So the columns of \(V\) are
\(\lambda_1-1\ge\dots\ge\lambda_q-1\ge s\ge \lambda_{q+1}-1\ge\dots\),
i.e.~\(V\) is the sorted diagram of \(B(\lambda)\) (columns of height 0
at the end are discarded). Finally, the cells moved in passing from
\(U\) to \(V\) are exactly the images under rule 2 of the cells
\((i,h)\) of \(\lambda\) with \(h-1\ge s+1\), i.e.~\(h\ge s+2\);
composing ``\((i,h)\mapsto(i+1,h-1)\) then one column left'' gives rule
3. \(\square\)

\subparagraph{\texorpdfstring{Lemma 2 (convergence to \(\delta_k\); used
only for context, not for the
bound)}{Lemma 2 (convergence to \textbackslash delta\_k; used only for context, not for the bound)}}\label{lemma-2-convergence-to-delta_k-used-only-for-context-not-for-the-bound}

Let \(W(\lambda)=\sum_{\text{cells}} d(i,h)\). By Lemma 1,
\(W(B(\lambda))=W(\lambda)-\#\{(i,h)\in\lambda: h\ge s+2\}\le W(\lambda)\).
Since diagonal \(d\) has only \(d\) slots, among all sets of \(T_k\)
cells with at most \(d\) cells on diagonal \(d\) the minimum of \(W\) is
attained only by \(\delta_k\) (fill diagonals \(1,\dots,k\)); hence
\(W(\lambda)\ge W(\delta_k)\) with equality iff \(\lambda=\delta_k\).
(Combined with the fact that a partition \(\ne\delta_k\) of \(T_k\)
cannot be periodic with all steps of cost 0 --- which is C1 material and
not needed below --- this gives convergence; we do not use it.)

\subparagraph{\texorpdfstring{Theorem A (lower bound). For every
\(k\ge1\), \(D_B(T_k)\ge k^2-k\), attained
by}{Theorem A (lower bound). For every k\textbackslash ge1, D\_B(T\_k)\textbackslash ge k\^{}2-k, attained by}}\label{theorem-a-lower-bound.-for-every-kge1-d_bt_kge-k2-k-attained-by}

\[\gamma_k=(k-1,\,k-1,\,k-2,\,k-3,\dots,2,\,1,\,1)\quad(k\ge2),\qquad \gamma_1=(1).\]
(\(\gamma_k\) is \(\delta_k\) with the part \(k\) replaced by the two
parts \(k-1\) and \(1\); \(|\gamma_k|=T_k\).)

\emph{Proof.} For \(k=1\) the only partition is \(\delta_1\) and
\(D_B(1)=0=k^2-k\). Let \(k\ge2\). For \(p_H\in\{1,\dots,k\}\) and
\(p_X\in\{1,\dots,k+1\}\) let \(S(p_H,p_X)\) be the set of cells of
\(\delta_k\) with the cell \((p_H,\,k+1-p_H)\) removed (a \emph{hole} in
slot \(p_H\) of diagonal \(k\)) and the cell \((p_X,\,k+2-p_X)\) added
(an \emph{extra cell} in slot \(p_X\) of diagonal \(k+1\)). Its column
heights are \(k+1-p\) for \(p\notin\{p_H,p_X\}\), \(k-p_H\) for
\(p=p_H\), \(k+2-p_X\) for \(p=p_X\) (and column \(k+1\) exists iff
\(p_X=k+1\)). It is a partition diagram iff \(p_X\ne p_H\) and
\(p_X\ne p_H+1\) (if \(p_X=p_H\) the cell \((p_H,k+2-p_H)\) would sit
above a missing cell; if \(p_X=p_H+1\) column \(p_X\) would be taller
than column \(p_H\); otherwise the heights are weakly decreasing because
column \(p_X\) has the same height as column \(p_X-1\) and column
\(p_H\) has the same height as column \(p_H+1\)). Note
\(\gamma_k=S(1,k+1)\): columns \(k-1,k-1,k-2,\dots,2,1,1\).

Number of columns of \(S(p_H,p_X)\): \(s=k+1\) if \(p_X=k+1\); \(s=k-1\)
if \(p_H=k\) (then \(p_X\le k\)); \(s=k\) otherwise. Heights: every
column has height \(\le k\) except column \(p_X\), of height
\(k+2-p_X\le k+1\), with equality iff \(p_X=1\). Hence a cell of height
\(\ge s+2\) exists only when \(s=k-1\) and \(p_X=1\), i.e.~only in
\(S(k,1)\), and then it is the single cell \((1,k+1)\).

\emph{Case \(S(k,1)\).} By Lemma 1 with \(s=k-1\): the bottom row
\((i,1)\), \(i\le k-1\), becomes column 1 with heights \(1,\dots,k-1\);
the cells \((i,h)\) with \(2\le h\le k\) (all remaining cells of
\(\delta_k\) minus the hole, since the hole is \((k,1)\) and column
\(k\) is empty) go to \((i+1,h-1)\), giving columns \(2,\dots,k\) of
heights \(k-1,\dots,1\); the cell \((1,k+1)\) drops to \((1,k)\),
completing column 1 to height \(k\). So \(B(S(k,1))=\delta_k\).

\emph{Other cases.} No cell has height \(\ge s+2\), so by Lemma 1 every
cell rotates one slot along its diagonal; the same is true of the empty
slot (hole) on diagonal \(k\), because the cells of diagonal \(k\)
permute cyclically among its slots. Hence \(B(S(p_H,p_X))=S(p_H',p_X')\)
with \(p_H'\equiv p_H+1\pmod k\), \(p_X'\equiv p_X+1\pmod{k+1}\)
(representatives in \(\{1..k\}\), \(\{1..k+1\}\)). (Lemma 1 guarantees
the result is a partition diagram, so the validity condition is
automatically preserved.)

Therefore, starting from \(\gamma_k=S(1,k+1)\), as long as no step has
hit \(S(k,1)\) we have \(B^t(\gamma_k)=S(p_H(t),p_X(t))\) with
\(p_H(t)\equiv 1+t\pmod k\), \(p_X(t)\equiv k+1+t\equiv t\pmod{k+1}\).
None of these equals \(\delta_k\) (they have a hole). The first
\(t\ge0\) with \(S(p_H(t),p_X(t))=S(k,1)\) satisfies
\(t\equiv k-1\pmod k\) and \(t\equiv1\pmod{k+1}\). Writing \(t=k-1+ak\):
\(k-1+ak\equiv -2-a\equiv 1\pmod{k+1}\), so
\(a\equiv-3\equiv k-2\pmod{k+1}\), and the least \(a\ge0\) is \(a=k-2\)
(valid since \(k\ge2\)). Thus \(t_0=k-1+(k-2)k=k^2-k-1\),
\(B^{t}(\gamma_k)\ne\delta_k\) for \(0\le t\le k^2-k-1\), and
\(B^{k^2-k}(\gamma_k)=B(S(k,1))=\delta_k\). Since \(\delta_k\) is a
fixed point and (by Lemma 2, or by C1) it is the only cyclic partition
of \(T_k\), \(d_B(\gamma_k)=k^2-k\). \(\square\)

(Consistency check with the statement's example and data:
\(\gamma_3=(2,2,1,1)\), \(d_B=6\); the exhaustive computation below
gives \(D_B(T_k)=k^2-k\) and confirms \(\gamma_k\) among the extremals
for \(k\le9\).)

\subparagraph{Theorem B (upper bound) --- NOT PROVED
HERE}\label{theorem-b-upper-bound-not-proved-here}

Claim: \(d_B(\lambda)\le k^2-k\) for every partition \(\lambda\) of
\(T_k\). This is the theorem of Igusa (1985) / Etienne (1991) (CITED,
proof not reproduced). What is established:

\begin{itemize}
\tightlist
\item
  \textbf{Exhaustive verification for \(k\le9\)} (exact integer
  arithmetic, all \(p(T_k)\) partitions, 0.5 s total):
  \(D_B(T_k)=k^2-k\) for \(k=1,\dots,9\). This proves the cell's
  statement only for \(k\le 9\).
\item
  Structural facts from Lemma 1 available for the next attempt: (i)
  \(W\) strictly decreases exactly at steps where some pile has size
  \(\ge s+2\), and \(\gamma_k\) has \(W(\gamma_k)-W(\delta_k)=1\), so
  the difficulty is bounding runs of ``free'' steps, not the number of
  costly steps; (ii) empirically (exhaustive, \(k\le8\)) the number of
  piles lies in \(\{k-1,k,k+1\}\) for all \(t\ge T_{k-1}-(k-2)\), which
  suggests the strategy: bound the time until the configuration becomes
  ``near-staircase'', then analyse the near-staircase dynamics as in
  Theorem A.
\end{itemize}

\subparagraph{Conclusion of this
attempt}\label{conclusion-of-this-attempt}

\(D_B(T_k)\ge k^2-k\) for all \(k\ge1\) (Theorem A, complete), with
explicit extremal \(\gamma_k\); \(D_B(T_k)=k^2-k\) for \(k\le9\) by
exhaustive computation; the general inequality \(D_B(T_k)\le k^2-k\)
remains to be proved.

\subparagraph{3. Verification: instructions, dependencies,
timings}\label{verification-instructions-dependencies-timings}

Python 3 standard library. Scripts (re-run by the orchestrator, see
\texttt{runs/p3\_c2/verifica/}): - code\_1 (python, rigor
\texttt{exact}): Exhaustive exact computation of D\_B(T\_k) and of all
extremal partitions for k ≤ 9 (all partitions of T\_k, orbit followed to
δ\_k). Finite set covered: every partition of T\_k for k=1..9 (up to
89134 partitions). Wall-clock 0.48 s. Confirms D\_B(T\_k)=k\^{}2−k for
k≤9; does not prove the general upper bound. - code\_2 (python, rigor
\texttt{exact}): Exploration only: along every orbit for k ≤ 8, records
the last time the pile count leaves \{k−1,k,k+1\} and the first time
diagonals 1..k−1 are full; used to guide the (unfinished) upper-bound
strategy. Exact integers, all partitions of T\_k for k=3..8,
\textasciitilde1 min.

\subparagraph{4. Sources and
contribution}\label{sources-and-contribution}

\begin{itemize}
\tightlist
\item
  B. Hopkins, 30 Years of Bulgarian Solitaire, College Math. J. 43
  (2012) 135--140 (read in full; p.137 statement that Igusa proved
  γ\_k=(k−1,k−1,k−2,\ldots,2,1,1) is at maximal distance k(k−1)) ---
  CITED for context only
\item
  K. Igusa, Solution of the Bulgarian solitaire conjecture, Math. Mag.
  58 (1985) 259--271 --- NOT read (paywalled); statement cited via
  Hopkins/Drensky
\item
  G. Etienne, Tableaux de Young et solitaire bulgare, J. Combin. Theory
  Ser. A 58 (1991) 181--197 --- NOT read
\item
  J. R. Griggs, C.-C. Ho, The cycling of partitions and compositions
  under repeated shifts, Adv. Appl. Math. 21 (1998) 205--227 --- NOT
  read
\item
  arXiv:1503.00885 (Drensky) and arXiv:2607.17194 (Meštrović) --- read;
  both only cite Igusa/Etienne for the k(k−1) bound and name Toom's
  extremal τ=γ\_k
\item
  M. Jonsson, Processes on Integer Partitions and Their Limit Shapes,
  PhD thesis, Mälardalen Univ. 2017 (DiVA diva2:1082060) --- read the
  relevant pages; cites Etienne for the game-tree height k\^{}2−k
\item
  Position with respect to the literature: The listed arXiv abstracts
  (math/0401385, 1703.07102: random variants; 1101.1546: Toom's
  convergence proof; 2208.14496: orbit growth; 1503.00885 and
  2607.17194: surveys) do not prove the bound. I fetched and read
  (pdftotext) Drensky 1503.00885, Meštrović 2607.17194, Hopkins ``30
  years of Bulgarian solitaire'' (College Math. J. 43 (2012) 135--140),
  Hopkins--Jones (EJC 13 (2006) R80), N. Pham's honors thesis and M.
  Jonsson's PhD thesis (DiVA 1082060). All state, CITED: Knuth
  conjectured and Igusa (Math. Mag. 58 (1985) 259--271) and Etienne
  (JCTA 58 (1991) 181--197) proved that for n=T\_k the maximal number of
  moves is k(k−1), attained by γ\_k=(k−1,k−1,k−2,\ldots,2,1,1) (Hopkins
  2012, p.~137: ``Igusa {[}15{]} shows that the partition γ\_k \ldots{}
  is at distance k(k−1) from τ\_k and that this distance is maximal'').
  None of the read sources reproduces the proof. My approach follows the
  classical ``cards on diagonals'' idea mentioned by Hopkins (p.~137)
  but makes it exact (Lemma 1) and uses it to prove the lower bound in
  full; the upper bound argument of Igusa/Etienne is not reproduced.
\end{itemize}

\subparagraph{5. Limits and unresolved
parts}\label{limits-and-unresolved-parts}

\begin{itemize}
\tightlist
\item
  Upper bound D\_B(T\_k) ≤ k\^{}2 − k for all k: not proved (only
  verified exhaustively for k ≤ 9 and cited from Igusa 1985 / Etienne
  1991, whose proofs I could not access). This is the next blocker.
\item
  Theorem A uses that δ\_k is the only cyclic partition of T\_k to
  conclude d\_B(γ\_k) = k\^{}2−k rather than merely
  B\textsuperscript{\{k}2−k\}(γ\_k)=δ\_k; this is the C1 statement,
  sketched via the potential W in Lemma 2 but not written out in full
  here (the strict-decrease/non-periodicity step is omitted).
\item
  Lemma 1's proof treats ties (columns of equal height) implicitly: the
  sorted diagram is determined by the multiset of column heights, so
  ties do not affect the cell set, but the reader may want this stated
  explicitly.
\end{itemize}

\subsubsection{Decisioni degli agenti}
\paragraph{Run \texttt{p3\_c2}}
\begin{description}
\item[Iterazione 1 — attempt\_001]
\textbf{Researcher}: famiglia \texttt{direct\_proof}, sotto-obiettivo: Establish the value D\_B(T\_k) = k\textasciicircum{}2 - k: prove the lower bound D\_B(T\_k) ≥ k\textasciicircum{}2 - k with an explicit extremal family and an exact cell-motion lemma for B; identify the general upper bound as the remaining blocker.. Stato dichiarato: \texttt{LEMMA\_CANDIDATE}.
\emph{Perché questa scelta}: First iteration, no blocker stated. The cell asks for the exact value with both bounds; the literature (read: Hopkins 2012 survey) fixes the answer k\textasciicircum{}2-k and Igusa's extremal γ\_k, but no accessible source contains a written proof (Igusa 1985, Etienne 1991, Griggs–Ho 1998 are paywalled; the surveys and theses I read only cite them). I therefore built the proof from scratch. The lower bound and an exact structural lemma are complete; the general upper bound is not, and I say so.
\emph{Rispetto allo stato dell'arte}: The listed arXiv abstracts (math/0401385, 1703.07102: random variants; 1101.1546: Toom's convergence proof; 2208.14496: orbit growth; 1503.00885 and 2607.17194: surveys) do not prove the bound. I fetched and read (pdftotext) Drensky 1503.00885, Meštrović 2607.17194, Hopkins "30 years of Bulgarian solitaire" (College Math. J. 43 (2012) 135–140), Hopkins–Jones (EJC 13 (2006) R80), N. Pham's honors thesis and M. Jonsson's PhD thesis (DiVA 1082060). All state, CITED: Knuth conjectured and Igusa (Math. Mag. 58 (1985) 259–271) and Etienne (JCTA 58 (1991) 181–197) proved that for n=T\_k the maximal number of moves is k(k−1), attained by γ\_k=(k−1,k−1,k−2,…,2,1,1) (Hopkins 2012, p. 137: "Igusa [15] shows that the partition γ\_k … is at distance k(k−1) from τ\_k and that this distance is maximal"). None of the read sources reproduces the proof. My approach follows the classical "cards on diagonals" idea mentioned by Hopkins (p. 137) but makes it exact (Lemma 1) and uses it to prove the lower bound in full; the upper bound argument of Igusa/Etienne is not reproduced.
\textbf{Referee}: verdetto \texttt{REJECT} (review\_status \texttt{NEEDS\_WORK}).
\emph{Errore fatale}: Theorem B (upper bound D\_B(T\_k) ≤ k\textasciicircum{}2−k for all k) is not proved; the submission itself marks it 'NOT PROVED HERE' and only cites Igusa/Etienne (statement not reproduced, sources not read) plus an exhaustive check for small k. The cell demands both bounds with proof, so the declared target 'main' (exact value k\textasciicircum{}2−k for every k) is not established. Finite exhaustive verification (k ≤ 8 in the orchestrator's re-run of code\_1; k ≤ 9 claimed) does not prove the universal upper bound. What would close it: a proof that every partition of T\_k reaches δ\_k within k\textasciicircum{}2−k shifts (e.g. via the potential W of Lemma 2 plus a bound on the length of runs of cost‑0 steps).
\emph{Prossimo passo richiesto}: Address the stated blocking obligation without silently changing the target
\end{description}

\subsubsection{Codice su cui poggia il tentativo}
\begin{lstlisting}[caption={attempt\_001 code\_1 [python, rigore exact] — Exhaustive exact computation of D\_B(T\_k) and of all extremal partitions for k ≤ 9 (all par — verdetto del Referee: REJECT}]
"""Calcolo esatto di D_B(T_k) per k piccoli (esplorazione, aritmetica intera esatta).

Enumera tutte le partizioni di n = T_k, applica lo shift B fino a raggiungere
delta_k (unico punto fisso per n triangolare: verifichiamo anche questo)
e registra la distanza massima e le partizioni che la raggiungono.
"""
import sys


def shift(parti):
    """Applica lo shift B: togli 1 da ogni parte, aggiungi la parte s."""
    s = len(parti)
    nuove = [p - 1 for p in parti if p > 1] + [s]
    return tuple(sorted(nuove, reverse=True))


def partizioni(n, massimo=None):
    """Genera tutte le partizioni di n con parti <= massimo (ordine decrescente)."""
    if massimo is None:
        massimo = n
    if n == 0:
        yield ()
        return
    for prima in range(min(n, massimo), 0, -1):
        for resto in partizioni(n - prima, prima):
            yield (prima,) + resto


def distanze_al_punto_fisso(n, punto_fisso):
    """Restituisce dict partizione -> numero di shift per arrivare al punto fisso.

    Usa memoizzazione lungo le orbite; presuppone che ogni orbita finisca in punto_fisso
    (lo verifichiamo con un controllo di sicurezza sul numero di passi).
    """
    dist = {punto_fisso: 0}
    for lam in partizioni(n):
        cammino = []
        cur = lam
        while cur not in dist:
            cammino.append(cur)
            cur = shift(cur)
            assert len(cammino) <= 10 * n * n, "orbita troppo lunga: non converge?"
        base = dist[cur]
        for i, p in enumerate(reversed(cammino)):
            dist[p] = base + i + 1
    return dist


def main(k_max):
    for k in range(1, k_max + 1):
        n = k * (k + 1) // 2
        delta = tuple(range(k, 0, -1))
        dist = distanze_al_punto_fisso(n, delta)
        D = max(dist.values())
        estremali = sorted(p for p, d in dist.items() if d == D)
        print(f"k={k} n={n} #partizioni={len(dist)} D_B={D} k^2-k={k*k-k} "
              f"#estremali={len(estremali)}")
        for p in estremali[:12]:
            print("   ", p)


if __name__ == "__main__":
    main(int(sys.argv[1]) if len(sys.argv) > 1 else 8)

\end{lstlisting}
\noindent\emph{Riesecuzione fidata dell'orchestratore}: \texttt{orchestrator re-ran code\_1.py (python3, clean copy of the researcher sandbox): exit 0 in 0.1s; stdout: '1, 1)\textbackslash{}n    (4, 4, 3, 3, 3, 3, 3, 3, 2)\textbackslash{}n    (4, 4, 4, 3, 2, 2, 2, 2, 2, 1, 1, 1)\textbackslash{}n    (4, 4, 4, 3, 2, 2, 2, 2, 2, 2, 1)\textbackslash{}nk=8 n=36 \#partizioni=17977 D\_B=56 k\textasciicircum{}2-k=56 \#estremali=1267\textbackslash{}n    (5, 4, 4, 3, 3, 3, 3, 2, 2, 2, 2, 1, 1, 1)\textbackslash{}n    (5, 4, 4, 3, 3, 3, 3, 2, 2, 2, 2, 2, 1)\textbackslash{}n    (5, 4, 4, 3, 3, 3,}
\begin{lstlisting}[caption={attempt\_001 code\_2 [python, rigore exact] — Exploration only: along every orbit for k ≤ 8, records the last time the pile count leaves — verdetto del Referee: REJECT}]
"""Esplorazione: statistiche sulla successione m(t) = numero di pile lungo le orbite."""
import sys
from collections import Counter
from calcola_DB_triangolari import shift, partizioni


def conteggio_diagonali(parti):
    """c_d = numero di celle (i,j) del diagramma (righe = pile) con i+j-1 = d."""
    c = Counter()
    for i, p in enumerate(parti, start=1):
        for j in range(1, p + 1):
            c[i + j - 1] += 1
    return c


def diagonali_piene_fino(parti, d_max):
    """True se le diagonali 1..d_max sono tutte piene (c_d = d)."""
    c = conteggio_diagonali(parti)
    return all(c[d] == d for d in range(1, d_max + 1))


def main(k):
    n = k * (k + 1) // 2
    delta = tuple(range(k, 0, -1))
    peggio_stab = 0; peggio_pieno = 0; combinazioni = Counter()
    for lam in partizioni(n):
        orbita = [lam]
        while orbita[-1] != delta:
            orbita.append(shift(orbita[-1]))
        d = len(orbita) - 1
        m = [len(p) for p in orbita]
        t_stab = max([t + 1 for t in range(len(m)) if abs(m[t] - k) > 1] + [0])
        t_pieno = next(t for t, p in enumerate(orbita) if diagonali_piene_fino(p, k - 1))
        peggio_stab = max(peggio_stab, t_stab)
        peggio_pieno = max(peggio_pieno, t_pieno)
        combinazioni[(t_pieno, d - t_pieno)] += 1
    print(f"k={k}: max t_stab={peggio_stab}, max t_pieno={peggio_pieno}, k^2-k={k*k-k}")


if __name__ == "__main__":
    for k in range(3, int(sys.argv[1]) + 1):
        main(k)

\end{lstlisting}
\noindent\emph{Riesecuzione fidata dell'orchestratore}: \texttt{orchestrator re-ran code\_2.py (python3, clean copy of the researcher sandbox): exit 1 in 0.1s; stdout: ''; stderr: 'Traceback (most recent call last):\textbackslash{}n  File "/Users/thomastumini/proof-pursuit/runs/p3\_c2/verifica/attempt\_001/code\_2.py", line 41, in \textless{}module\textgreater{}\textbackslash{}n    for k in range(3, int(sys.argv[1]) + 1):\textbackslash{}n                          \textasciitilde{}\textasciitilde{}\textasciitilde{}\textasciitilde{}\textasciitilde{}\textasciitilde{}\textasciitilde{}\textasciitilde{}\textasciicircum{}\textasciicircum{}\textasciicircum{}\textbackslash{}nIndexError: list index out of range'}

\section{arXiv literature consulted}
\begin{thebibliography}{99}
\bibitem{math/0401385v2} Serguei Popov. \emph{Random Bulgarian solitaire}. arXiv:math/0401385v2 (2004). \url{http://arxiv.org/abs/math/0401385v2}
\bibitem{1503.00885v1} Vesselin Drensky. \emph{The Bulgarian solitaire and the mathematics around it}. arXiv:1503.00885v1 (2015). \url{http://arxiv.org/abs/1503.00885v1}
\bibitem{2607.17194v1} Romeo Meštrović. \emph{A short survey the game Bulgarian solitaire and related games}. arXiv:2607.17194v1 (2026). \url{http://arxiv.org/abs/2607.17194v1}
\bibitem{1101.1546v3} Therese A. Hart, Gabriel Khan, Mizan R. Khan. \emph{Revisiting Toom's proof of Bulgarian Solitaire}. arXiv:1101.1546v3 (2011). \url{http://arxiv.org/abs/1101.1546v3}
\bibitem{1703.07102v1} Kimmo Eriksson, Markus Jonsson abd Jonas Sjöstrand. \emph{An exponential limit shape of random \$q\$-proportion Bulgarian solitaire}. arXiv:1703.07102v1 (2017). \url{http://arxiv.org/abs/1703.07102v1}
\bibitem{2208.14496v1} Nhung Pham. \emph{Limiting behavior in growth of Bulgarian Solitaire orbits}. arXiv:2208.14496v1 (2022). \url{http://arxiv.org/abs/2208.14496v1}
\end{thebibliography}
\end{document}
